\input{preamble}

\title{Section 3.5 --- The Quadratic Formula}

\begin{document}
\maketitle
\section{Using the Quadratic Formula}
Use the quadratic formula
\[
x=\frac{-b\pm\sqrt{b^2-4ac}}{2a}
\]
to solve each equation.  Simplify your answer as much as possible, but leave in exact radicals.
\begin{smenum}
\mitemxxxx
{$x^2-8x+12=0$}
{$6x^2+17x-14=0$}
{$2x^2+9x+3=0$}
{$5x^2-6x+3=0$}
\end{smenum}

\section{The Quadratic Formula and Standard Form}
Use the quadratic formula to solve each equation.  You will first need to write the equation in standard form.  Simplify your answer as much as possible, but leave in exact radicals.
\begin{smenum}
\mitemxxxx
{$3-x^2=0$}
{$x-x^2-5=0$}
{$x^2-3x+2=6$}
{$(7-x)(5-x)=24$}
\mitemxx
{$2x^2-15x-7=x^2-19x+3$}
{$(x+4)^2=2(x+4)(x-4)$}
\end{smenum}

\section{Using the Discriminant}
For each quadratic equation:
\begin{clist}
\item Compute the discriminant, $b^2-4ac$
\item Use the discriminant to determine whether the equation will have zero, one, or two real solutions.
\end{clist}
\begin{smenum}
\mitemxxxx
{$x^2-3x-2=0$}
{$-4x^2-3x-7=0$}
{$x^2-10x+25=0$}
{$-8x^2+7x-1=0$}
\end{smenum}

\section{Using the Quadratic Formula}
\begin{clist}
\item If it isn't already, write each equation in standard form.
\item Use the quadratic formula to solve.  Give your solutions as decimal approximations rounded to the nearest hundredth, if necessary.
\item Plug your answers into the original equation to show that they are correct.
\end{clist}
\begin{smenum}
\mitemxx
{$0.003x^2-28x+45=0$}
{$-4.9x^2+35x+3=0$}
\mitemxx
{$1.02x^2-2.54x-11.28=0$}
{$3x^2-48x=11x-19$}
\end{smenum}

\end{document}